\documentclass[11pt,paper=a4,answers]{exam}
\usepackage{graphicx,lastpage}
\usepackage{upgreek}
\usepackage{censor}
\usepackage{gensymb}
\usepackage{amsmath}
\usepackage{multicol}
\usepackage{enumitem}
\usepackage{tasks}
\usepackage{adjustbox}
\usepackage{xcolor}
\usepackage{tfrupee}
\setlength{\voffset}{0.25in}
\usepackage{tabularx}
\newcolumntype{Y}{>{\centering\arraybackslash}X}
%==============================================================
% Customise Non-Essential Packages Below
% =============================================================
\usepackage[most]{tcolorbox}
\usepackage{eso-pic}
\usepackage{tikz}
\AddToShipoutPictureBG{%
\begin{tikzpicture}[remember picture,overlay]
\foreach \x in {-8,-4,0,4,8,12,16,20}
\foreach \y in {-8,-4,0,4,8,12,16,20} {
\node[opacity=0.05,rotate=45,anchor=center] at ([xshift=\x cm,yshift=\y cm]current page.center) {%
\begin{minipage}{2cm}
\centering
\includegraphics[width=2cm]{images/logo_black.png}\\
\small preppeo.com
\end{minipage}%
};
}
\end{tikzpicture}%
}
\printanswers

%==============================================================
% Customise Non-Essential Packages Above
% =============================================================
\definecolor{svgray}{rgb}{0.90, 0.90, 0.90}
\graphicspath{{images/}}

\usepackage[normalem]{ulem}
\renewcommand\ULthickness{1pt}   %%---> For changing thickness of underline
\setlength\ULdepth{3.5ex}%\maxdimen ---> For changing depth of underline
\renewcommand{\baselinestretch}{1}
\chead{
\includegraphics[scale=0.1,height=2cm ]{logo.png}
}
\pagestyle{empty}

\pagestyle{headandfoot}
\headrule
\cfoot{W: https://preppeo.com; M: +91-8130711689}

\begin{document}
%==============================================================
% Customise Question paper details below this
% =============================================================
\begin{center}
    \centering
    \renewcommand{\arraystretch}{1.5}
    \begin{tabularx}{\textwidth}{YYY}
    & \normalsize\bf{{\Large{{Quantitative Reasoning}}}} & \\
    By Preppeo & \normalsize\bf{{sat\_lid\_029}} & SAT\\
    \normalsize{{\textbf{{Unit:}} Advanced Math}} & \normalsize{{\textbf{{Chapter:}} Function Notation	Function notation and interpreting functions}} & \normalsize{{\textbf{{Topic:}} Domain \& Range}} \\
    \end{tabularx}
\end{center}
\par\noindent\rule{\textwidth}{1pt}
% =============================================================
% Customise Question paper details above this
% =============================================================

\begin{questions}
\pointsinrightmargin
\pointsdroppedatright
\marksnotpoints
\pointpoints{mark}{marks}
\pointformat{\boldmath\themarginpoints}
\bracketedpoints

% =============================================================
% Question Begin
% =============================================================
% Lid Topic: Advanced Math — Domain & Range
% Total Questions: 50

% --- PART 1: Foundations and Algebraic Constraints (1-25) ---

% Question 1
% Category: Concept | Difficulty: Medium | Question Type: Multiple Choice
\question
The function $f$ is defined by $f(x) = \frac{1}{x - 9}$. Which of the following is the domain of $f$?
\begin{tasks}(1)
    \task All real numbers
    \task All real numbers such that $x \neq 0$
    \task All real numbers such that $x \neq 9$
    \task All real numbers such that $x > 9$
\end{tasks}

\begin{solution}
\textbf{Conceptual Explanation:} \\
The domain of a function consists of all possible input values ($x$) for which the function is defined. For rational functions (fractions with variables in the denominator), the function is undefined wherever the denominator equals zero.

\vspace{30pt}

\textbf{Calculation and Logic:} \\
Set the denominator equal to zero to find the excluded value: \\
$x - 9 = 0$ \\
$x = 9$

\vspace{15pt}

Since division by zero is undefined, $x$ can be any real number except 9.

\vspace{30pt}

Answer: \colorbox{yellow!100}{C}
\end{solution}

\vspace{40pt}

% Question 2
% Category: Exam level | Difficulty: Hard | Question Type: Numeric Entry
\question
$f(x) = \sqrt{2x - 14}$ \\
For the function $f$ defined above, what is the smallest value in the domain?

\begin{solution}
\textbf{Conceptual Explanation:} \\
For a square root function to result in a real number, the expression inside the radical (the radicand) must be greater than or equal to zero. 

\vspace{30pt}

\textbf{Calculation and Logic:} \\
Set up the inequality for the domain: \\
$2x - 14 \geq 0$

\vspace{15pt}

Solve for $x$: \\
$2x \geq 14$ \\
$x \geq 7$

\vspace{15pt}

The domain consists of all values 7 or greater. The smallest value is 7.

\vspace{30pt}

Answer: \colorbox{yellow!100}{7}
\end{solution}

\vspace{40pt}

% Question 3
% Category: Practice | Difficulty: Medium | Question Type: Multiple Choice
\question
\begin{center}
\begin{tikzpicture}[scale=0.7]
    \draw[->] (-1,0) -- (5,0) node[right] {$x$};
    \draw[->] (0,-1) -- (0,5) node[above] {$y$};
    \draw[ultra thick, blue] (1,2) -- (4,4);
    \filldraw[blue] (1,2) circle (2pt);
    \filldraw[blue] (4,4) circle (2pt);
\end{tikzpicture}
\end{center}
The graph of the function $f$ is shown in the $xy$-plane above. What is the range of $f$?
\begin{tasks}(2)
    \task $1 \leq x \leq 4$
    \task $2 \leq y \leq 4$
    \task $1 \leq y \leq 4$
    \task $2 \leq x \leq 4$
\end{tasks}

\begin{solution}
\textbf{Conceptual Explanation:} \\
The **Range** of a function is the set of all possible output values ($y$-coordinates) shown on the graph. We look at the vertical span of the graph from the lowest point to the highest point.

\vspace{30pt}

\textbf{Calculation and Logic:} \\
1. **Identify the lowest point:** The segment starts at $(1, 2)$. The $y$-value is 2. \\
2. **Identify the highest point:** The segment ends at $(4, 4)$. The $y$-value is 4.

\vspace{15pt}

The graph covers all $y$-values from 2 to 4, inclusive.

\vspace{30pt}

Answer: \colorbox{yellow!100}{B}
\end{solution}

\vspace{40pt}

% Question 4
% Category: Exam level | Difficulty: Hard | Question Type: Numeric Entry
\question
$f(x) = |x - 5| + 12$ \\
What is the minimum value of the range of the function $f$?

\begin{solution}
\textbf{Conceptual Explanation:} \\
An absolute value expression $|x - h|$ has a minimum value of 0. To find the minimum of the range, we determine the smallest possible value for the entire function.

\vspace{30pt}

\textbf{Calculation and Logic:} \\
The smallest value $|x - 5|$ can achieve is 0 (which happens when $x = 5$). \\
$f(5) = 0 + 12 = 12$.

\vspace{15pt}

For any other value of $x$, $|x - 5|$ will be positive, making $f(x)$ greater than 12. \\
Therefore, the range is $y \geq 12$, and the minimum value is 12.

\vspace{30pt}

Answer: \colorbox{yellow!100}{12}
\end{solution}

\vspace{40pt}

% Question 5
% Category: Concept | Difficulty: Medium | Question Type: Multiple Choice
\question
Which of the following values of $x$ is NOT in the domain of the function $f(x) = \frac{x+5}{x^2 - 16}$?
\begin{tasks}(4)
    \task -5
    \task 0
    \task 4
    \task 5
\end{tasks}

\begin{solution}
\textbf{Calculation and Logic:} \\
A rational function is undefined when the denominator is zero. \\
$x^2 - 16 = 0$ \\
$x^2 = 16$ \\
$x = 4$ or $x = -4$.

\vspace{15pt}

From the given options, 4 is an excluded value.

\vspace{30pt}

Answer: \colorbox{yellow!100}{C}
\end{solution}

\vspace{40pt}

% Question 6
% Category: Exam level | Difficulty: Hard | Question Type: Numeric Entry
\question
$g(x) = \frac{1}{(x-3)(x-10)}$ \\
The function $g$ is undefined for $x = a$ and $x = b$. What is the value of $a + b$?

\begin{solution}
\textbf{Calculation and Logic:} \\
The function is undefined where the denominator equals zero. \\
$(x - 3)(x - 10) = 0$ \\
$x = 3$ and $x = 10$.

\vspace{15pt}

Let $a = 3$ and $b = 10$. \\
$a + b = 3 + 10 = 13$.

\vspace{30pt}

Answer: \colorbox{yellow!100}{13}
\end{solution}

\vspace{40pt}

% Question 7
% Category: Practice | Difficulty: Medium | Question Type: Multiple Choice
\question
$f(x) = x^2 + 8$ \\
Which of the following represents the range of the function $f$?
\begin{tasks}(2)
    \task $y \geq 0$
    \task $y \geq 8$
    \task $x \geq 0$
    \task All real numbers
\end{tasks}

\begin{solution}
\textbf{Conceptual Explanation:} \\
The square of any real number ($x^2$) is always greater than or equal to zero. Adding a constant shifts the entire set of output values.

\vspace{30pt}

\textbf{Calculation and Logic:} \\
Minimum of $x^2$ is 0. \\
Minimum of $x^2 + 8$ is $0 + 8 = 8$. \\
The function can produce any value from 8 upwards.

\vspace{30pt}

Answer: \colorbox{yellow!100}{B}
\end{solution}

\vspace{40pt}

% Question 8
% Category: Exam level | Difficulty: Hard | Question Type: Numeric Entry
\question
A function $f$ has domain $2 \leq x \leq 8$. If $g(x) = f(x - 5)$, what is the maximum value in the domain of $g$?

\begin{solution}
\textbf{Conceptual Explanation:} \\
Transformations inside the function ($x - h$) shift the domain. For $f(x-5)$, the input seen by the original function $f$ must still fall between 2 and 8.

\vspace{30pt}

\textbf{Calculation and Logic:} \\
The input to $f$ is $(x - 5)$. \\
$2 \leq x - 5 \leq 8$

\vspace{15pt}

Add 5 to all parts of the inequality: \\
$2 + 5 \leq x \leq 8 + 5$ \\
$7 \leq x \leq 13$

\vspace{15pt}

The new domain is $[7, 13]$. The maximum value is 13.

\vspace{30pt}

Answer: \colorbox{yellow!100}{13}
\end{solution}

\vspace{40pt}

% Question 9
% Category: Concept | Difficulty: Medium | Question Type: Multiple Choice
\question
The function $h$ is defined by $h(x) = \sqrt{x}$. If $h(x)$ must be a real number, what is the domain of $h$?
\begin{tasks}(4)
    \task $x < 0$
    \task $x > 0$
    \task $x \geq 0$
    \task All real numbers
\end{tasks}

\begin{solution}
\textbf{Conceptual Explanation:} \\
The square root of a negative number is not a real number. Therefore, the input $x$ must be zero or positive.

\vspace{30pt}

Answer: \colorbox{yellow!100}{C}
\end{solution}

\vspace{40pt}

% Question 10
% Category: Exam level | Difficulty: Hard | Question Type: Numeric Entry
\question
$f(x) = -(x-4)^2 + 10$ \\
What is the maximum value in the range of function $f$?

\begin{solution}
\textbf{Conceptual Explanation:} \\
This is a downward-opening parabola (indicated by the negative sign before the squared term). The maximum value is the $y$-coordinate of the vertex.

\vspace{30pt}

\textbf{Calculation and Logic:} \\
Vertex form: $a(x-h)^2 + k$. \\
Vertex is $(4, 10)$. \\
Since the parabola opens downward, the highest point is 10. All other $y$-values are less than 10.

\vspace{30pt}

Answer: \colorbox{yellow!100}{10}
\end{solution}

% Question 11
% Category: Practice | Difficulty: Medium | Question Type: Multiple Choice
\question
\begin{center}
\begin{tikzpicture}[scale=0.6]
    \draw[step=1cm,gray!20,very thin] (-3,-3) grid (3,3);
    \draw[->] (-3,0) -- (3,0) node[right] {$x$};
    \draw[->] (0,-3) -- (0,3) node[above] {$y$};
    \draw[ultra thick, red] (-2,1) circle (1.5);
\end{tikzpicture}
\end{center}
The graph above shows a circle. Is this circle a representation of a function $y = f(x)$?
\begin{tasks}(1)
    \task Yes, because every $x$ has a $y$.
    \task No, because it fails the vertical line test.
    \task Yes, because it is a closed shape.
    \task No, because the domain is limited.
\end{tasks}

\begin{solution}
\textbf{Conceptual Explanation:} \\
A relation is a function if and only if every input $x$ corresponds to exactly one output $y$. Visually, this is checked using the vertical line test. If any vertical line crosses the graph more than once, it is not a function.

\vspace{30pt}

Answer: \colorbox{yellow!100}{B}
\end{solution}

\vspace{40pt}

% Question 12
% Category: Exam level | Difficulty: Hard | Question Type: Numeric Entry
\question
$f(x) = \frac{x-1}{x^2 - kx + 9}$ \\
If the function $f$ is undefined for exactly one value of $x$, and $k > 0$, what is the value of $k$?

\begin{solution}
\textbf{Conceptual Explanation:} \\
A function is undefined where the denominator equals zero. For a quadratic denominator to equal zero at exactly one point, it must be a perfect square trinomial (its discriminant must be zero).

\vspace{30pt}

\textbf{Calculation and Logic:} \\
Denominator: $x^2 - kx + 9$. \\
For a perfect square $(x-3)^2 = x^2 - 6x + 9$. \\
Or, using discriminant $D = b^2 - 4ac$: \\
$(-k)^2 - 4(1)(9) = 0$ \\
$k^2 - 36 = 0 \Rightarrow k = 6$ (since $k > 0$).

\vspace{30pt}

Answer: \colorbox{yellow!100}{6}
\end{solution}

\vspace{40pt}

% Question 13
% Category: Practice | Difficulty: Medium | Question Type: Multiple Choice
\question
$f(x) = 10x + 5$ \\
If the domain of $f$ is restricted to $0 \leq x \leq 3$, what is the range of $f$?
\begin{tasks}(2)
    \task $0 \leq y \leq 3$
    \task $5 \leq y \leq 35$
    \task $0 \leq y \leq 35$
    \task $10 \leq y \leq 30$
\end{tasks}

\begin{solution}
\textbf{Calculation and Logic:} \\
For a linear function with a positive slope, the minimum $y$ occurs at the minimum $x$, and the maximum $y$ occurs at the maximum $x$. \\
1. $f(0) = 10(0) + 5 = 5$. \\
2. $f(3) = 10(3) + 5 = 35$.

\vspace{15pt}

Range: $5 \leq y \leq 35$.

\vspace{30pt}

Answer: \colorbox{yellow!100}{B}
\end{solution}

\vspace{40pt}

% Question 14
% Category: Exam level | Difficulty: Hard | Question Type: Numeric Entry
\question
$h(x) = \sqrt{x+4} - 2$ \\
What is the $y$-intercept of the function $h$?

\begin{solution}
\textbf{Conceptual Explanation:} \\
The $y$-intercept is the output value when the input $x$ is 0. This is always a single point in the range.

\vspace{30pt}

\textbf{Calculation and Logic:} \\
$h(0) = \sqrt{0 + 4} - 2$ \\
$h(0) = \sqrt{4} - 2$ \\
$h(0) = 2 - 2 = 0$.

\vspace{30pt}

Answer: \colorbox{yellow!100}{0}
\end{solution}

\vspace{40pt}

% Question 15
% Category: Practice | Difficulty: Hard | Question Type: Multiple Choice
\question
Which of the following functions has a range of all real numbers?
\begin{tasks}(2)
    \task $f(x) = x^2$
    \task $f(x) = |x|$
    \task $f(x) = \sqrt{x}$
    \task $f(x) = x^3$
\end{tasks}

\begin{solution}
\textbf{Conceptual Explanation:} \\
Even powers (like $x^2$) and absolute values always produce non-negative results, limiting their range. Odd powers (like $x^3$) can produce any positive or negative value.

\vspace{30pt}

Answer: \colorbox{yellow!100}{D}
\end{solution}

\vspace{40pt}

% Question 16
% Category: Exam level | Difficulty: Hard | Question Type: Numeric Entry
\question
$f(x) = x^2 - 6x + 14$ \\
What is the minimum value in the range of $f$?

\begin{solution}
\textbf{Calculation and Logic:} \\
Find the vertex $x$-coordinate: $x = -b/2a = 6/2 = 3$. \\
Substitute $x=3$ into the function: \\
$f(3) = 3^2 - 6(3) + 14$ \\
$f(3) = 9 - 18 + 14 = 5$.

\vspace{15pt}

Since the parabola opens upward ($a=1 > 0$), the minimum value is 5.

\vspace{30pt}

Answer: \colorbox{yellow!100}{5}
\end{solution}

\vspace{40pt}

% Question 17
% Category: Practice | Difficulty: Medium | Question Type: Multiple Choice
\question
For which of the following values of $x$ is the function $f(x) = \sqrt{10-x}$ undefined?
\begin{tasks}(4)
    \task 0
    \task 5
    \task 10
    \task 15
\end{tasks}

\begin{solution}
\textbf{Calculation and Logic:} \\
Function is defined if $10 - x \geq 0 \Rightarrow 10 \geq x$. \\
The function is undefined for any $x > 10$. \\
Among the options, 15 is greater than 10.

\vspace{30pt}

Answer: \colorbox{yellow!100}{D}
\end{solution}

\vspace{40pt}

% Question 18
% Category: Exam level | Difficulty: Hard | Question Type: Numeric Entry
\question
If the range of $f(x) = 2x + b$ is $10 \leq y \leq 20$ for the domain $0 \leq x \leq 5$, what is the value of $b$?

\begin{solution}
\textbf{Calculation and Logic:} \\
Use the lower bound of the domain to find the lower bound of the range. \\
$f(0) = 2(0) + b = 10$ \\
$b = 10$.

\vspace{15pt}

(Verify with upper bound: $f(5) = 2(5) + 10 = 10 + 10 = 20$. Correct).

\vspace{30pt}

Answer: \colorbox{yellow!100}{10}
\end{solution}

\vspace{40pt}

% Question 19
% Category: Concept | Difficulty: Medium | Question Type: Multiple Choice
\question
\begin{center}
\begin{tikzpicture}[scale=0.6]
    \draw[->] (-3,0) -- (3,0) node[right] {$x$};
    \draw[->] (0,-3) -- (0,3) node[above] {$y$};
    \draw[ultra thick, blue] (-2,-2) -- (2,2);
    \node at (1,0.5) {$f$};
\end{tikzpicture}
\end{center}
Based on the graph, which of the following is true about the domain and range of $f$?
\begin{tasks}(1)
    \task The domain is the same as the range.
    \task The domain is larger than the range.
    \task The range is larger than the domain.
    \task Neither domain nor range can be determined.
\end{tasks}

\begin{solution}
\textbf{Calculation and Logic:} \\
Domain span (x): -2 to 2 (width = 4). \\
Range span (y): -2 to 2 (height = 4). \\
They are identical.

\vspace{30pt}

Answer: \colorbox{yellow!100}{A}
\end{solution}

\vspace{40pt}

% Question 20
% Category: Exam level | Difficulty: Hard | Question Type: Numeric Entry
\question
$f(x) = \frac{x+2}{x^2-4}$ \\
For what value of $x$ is the function $f$ undefined, even though the numerator is zero at that point?

\begin{solution}
\textbf{Conceptual Explanation:} \\
A function is undefined if the denominator is zero, regardless of what the numerator is. If both are zero, it is still undefined (this often indicates a "hole" in the graph).

\vspace{30pt}

\textbf{Calculation and Logic:} \\
Numerator is zero at $x = -2$. \\
Denominator $x^2 - 4 = 0$ at $x = 2$ and $x = -2$. \\
Since $x = -2$ makes the denominator zero, the function is undefined there.

\vspace{30pt}

Answer: \colorbox{yellow!100}{-2}
\end{solution}

\vspace{40pt}

% Question 21
% Category: Practice | Difficulty: Medium | Question Type: Multiple Choice
\question
If $f(x) = x^2 - 4$, what is the range of $f$ if the domain is all real numbers?
\begin{tasks}(4)
    \task $y \geq 0$
    \task $y \geq -4$
    \task $y \leq -4$
    \task All real numbers
\end{tasks}

\begin{solution}
\textbf{Calculation and Logic:} \\
Min of $x^2 = 0$. \\
Min of $x^2 - 4 = -4$. \\
Range: $y \geq -4$.

\vspace{30pt}

Answer: \colorbox{yellow!100}{B}
\end{solution}

\vspace{40pt}

% Question 22
% Category: Exam level | Difficulty: Hard | Question Type: Numeric Entry
\question
If $f(x) = \sqrt{x-a}$ has a domain of $x \geq 15$, what is the value of $a$?

\begin{solution}
\textbf{Calculation and Logic:} \\
$x - a \geq 0 \Rightarrow x \geq a$. \\
Since the domain is $x \geq 15$, $a$ must be 15.

\vspace{30pt}

Answer: \colorbox{yellow!100}{15}
\end{solution}

\vspace{40pt}

% Question 23
% Category: Practice | Difficulty: Medium | Question Type: Multiple Choice
\question
$f(x) = \frac{1}{\sqrt{x-2}}$ \\
What is the domain of $f$?
\begin{tasks}(4)
    \task $x \geq 2$
    \task $x > 2$
    \task $x \neq 2$
    \task All real numbers
\end{tasks}

\begin{solution}
\textbf{Conceptual Explanation:} \\
The value under the radical must be $\geq 0$. However, since the radical is in the denominator, the total expression cannot be zero. Thus, $x-2$ must be strictly greater than 0.

\vspace{30pt}

Answer: \colorbox{yellow!100}{B}
\end{solution}

\vspace{40pt}

% Question 24
% Category: Exam level | Difficulty: Hard | Question Type: Numeric Entry
\question
$f(x) = 5^x$ \\
What is the smallest integer that is NOT in the range of function $f$?

\begin{solution}
\textbf{Conceptual Explanation:} \\
Exponential functions in the form $a^x$ (where $a > 0$) always produce positive results. They can get very close to 0 but never reach it or become negative.

\vspace{30pt}

\textbf{Calculation and Logic:} \\
Range: $y > 0$. \\
The values 1, 2, 3... are in the range. \\
The value 0 is the largest integer not in the range (since everything $\leq 0$ is excluded).

\vspace{30pt}

Answer: \colorbox{yellow!100}{0}
\end{solution}

\vspace{40pt}

% Question 25
% Category: Concept | Difficulty: Medium | Question Type: Multiple Choice
\question
If the graph of a function is a horizontal line $y = 3$, what is its range?
\begin{tasks}(4)
    \task All real numbers
    \task $y \geq 3$
    \task $\{3\}$
    \task $x = 3$
\end{tasks}

\begin{solution}
\textbf{Conceptual Explanation:} \\
The range consists only of the output values. Since every input produces the same output (3), the range contains only that single number.

\vspace{30pt}

Answer: \colorbox{yellow!100}{C}
\end{solution}

% --- PART 2: Advanced Graphical and Modeling Analysis (26-50) ---

% Question 26
% Category: Exam level | Difficulty: Hard | Question Type: Numeric Entry
\question
The function $f$ is defined by $f(x) = \sqrt{40 - 5x}$. What is the largest integer value of $x$ that is in the domain of $f$?

\begin{solution}
\textbf{Conceptual Explanation:} \\
For a square root function to be defined in the set of real numbers, the expression under the radical (the radicand) must be greater than or equal to zero. 

\vspace{30pt}

\textbf{Calculation and Logic:} \\
Set up the inequality for the domain: \\
$40 - 5x \geq 0$

\vspace{15pt}

Subtract 40 from both sides: \\
$-5x \geq -40$

\vspace{15pt}

Divide by $-5$. \textbf{Remember:} When dividing an inequality by a negative number, the inequality sign must be flipped: \\
$x \leq 8$

\vspace{15pt}

The domain consists of all values less than or equal to 8. Therefore, the largest integer value in the domain is 8.

\vspace{30pt}

Answer: \colorbox{yellow!100}{8}
\end{solution}

\vspace{40pt}

% Question 27
% Category: Practice | Difficulty: Medium | Question Type: Multiple Choice
\question
\begin{center}
\begin{tikzpicture}[scale=0.6]
    \draw[->] (-1,0) -- (6,0) node[right] {$x$};
    \draw[->] (0,-1) -- (0,5) node[above] {$y$};
    \draw[ultra thick, blue] (1,4) .. controls (2,0) and (4,0) .. (5,4);
    \filldraw[blue] (1,4) circle (2pt);
    \filldraw[blue] (5,4) circle (2pt);
    \draw[dashed] (3,1) -- (0,1) node[left] {1};
\end{tikzpicture}
\end{center}
The graph of the function $f$ is shown in the $xy$-plane above. What is the range of $f$?
\begin{tasks}(2)
    \task $1 \leq y \leq 4$
    \task $1 \leq x \leq 5$
    \task $0 \leq y \leq 4$
    \task $4 \leq y \leq 5$
\end{tasks}

\begin{solution}
\textbf{Conceptual Explanation:} \\
The range represents the vertical span of the function. Look for the absolute lowest $y$-value reached by the curve and the absolute highest $y$-value.

\vspace{30pt}

\textbf{Calculation and Logic:} \\
1. \textbf{Identify the lowest point:} According to the dashed line and the curve, the minimum $y$-value reached is 1. \\
2. \textbf{Identify the highest point:} The graph ends at $y = 4$ on both sides.

\vspace{15pt}

The function outputs all values from 1 to 4 inclusive.

\vspace{30pt}

Answer: \colorbox{yellow!100}{A}
\end{solution}

\vspace{40pt}

% Question 28
% Category: Exam level | Difficulty: Hard | Question Type: Numeric Entry
\question
$f(x) = \frac{10}{x^2 - 14x + 49}$ \\
The function $f$ is undefined for only one value of $x$. What is that value?

\begin{solution}
\textbf{Calculation and Logic:} \\
The function is undefined when the denominator is zero. \\
$x^2 - 14x + 49 = 0$

\vspace{15pt}

Factor the quadratic expression: \\
$(x - 7)(x - 7) = 0$ \\
$(x - 7)^2 = 0$

\vspace{15pt}

Solving for $x$ gives: \\
$x = 7$

\vspace{30pt}

Answer: \colorbox{yellow!100}{7}
\end{solution}

\vspace{40pt}

% Question 31 (Example of shifted numbering for variety)
% Category: Concept | Difficulty: Medium | Question Type: Multiple Choice
\question
\begin{center}
\begin{tikzpicture}[scale=0.6]
    \draw[->] (-4,0) -- (4,0) node[right] {$x$};
    \draw[->] (0,-1) -- (0,5) node[above] {$y$};
    \draw[ultra thick, red] (-3,4) -- (0,1) -- (3,4);
    \node at (1,3) {$g(x)$};
\end{tikzpicture}
\end{center}
The graph of $g(x) = |x| + c$ is shown. Based on the graph, what is the range of $g$?
\begin{tasks}(4)
    \task $y \geq 4$
    \task $y \leq 1$
    \task $y \geq 1$
    \task All real numbers
\end{tasks}

\begin{solution}
\textbf{Conceptual Explanation:} \\
The vertex of this absolute value function is $(0, 1)$. Since the "V" shape opens upward, the function reaches a minimum value of 1 and continues toward infinity.

\vspace{30pt}

Answer: \colorbox{yellow!100}{C}
\end{solution}

\vspace{40pt}

% Question 30
% Category: Exam level | Difficulty: Hard | Question Type: Numeric Entry
\question
A scientist uses the function $H(t) = -4.9t^2 + 20t + 5$ to model the height of a projectile in meters $t$ seconds after it is launched. If the projectile hits the ground at $t \approx 4.3$ seconds, what is the maximum integer in the domain of $H$ that represents a physically possible time?

\begin{solution}
\textbf{Conceptual Explanation:} \\
In real-world word problems, the domain is often restricted by the context. Time cannot be negative, and the model stops being valid once the object hits the ground ($H=0$).

\vspace{30pt}

\textbf{Calculation and Logic:} \\
Physical domain: $0 \leq t \leq 4.3$. \\
The integers in this range are $\{0, 1, 2, 3, 4\}$. \\
The largest integer is 4.

\vspace{30pt}

Answer: \colorbox{yellow!100}{4}
\end{solution}

\vspace{40pt}

% Question 31
% Category: Practice | Difficulty: Hard | Question Type: Multiple Choice
\question
$f(x) = \frac{x^2 - 1}{x - 1}$ \\
Which of the following is true about the domain of $f$?
\begin{tasks}(1)
    \task The domain is all real numbers.
    \task The domain is all real numbers except $x=1$.
    \task The domain is all real numbers except $x=1$ and $x=-1$.
    \task The domain is $x > 1$.
\end{tasks}

\begin{solution}
\textbf{Conceptual Explanation:} \\
Even if a rational expression can be simplified (e.g., by factoring and canceling), the domain is \textbf{always} determined by the \textbf{original} denominator before any simplification occurs.

\vspace{30pt}

\textbf{Calculation and Logic:} \\
Original denominator: $x - 1$. \\
$x - 1 = 0 \Rightarrow x = 1$. \\
The value $x=1$ must be excluded.

\vspace{30pt}

Answer: \colorbox{yellow!100}{B}
\end{solution}

\vspace{40pt}

% Question 32
% Category: Exam level | Difficulty: Hard | Question Type: Numeric Entry
\question
$f(x) = \sqrt{x - 3}$ \\
$g(x) = f(x + 10)$ \\
What is the smallest value in the domain of function $g$?

\begin{solution}
\textbf{Calculation and Logic:} \\
Step 1: Write the expression for $g(x)$: \\
$g(x) = \sqrt{(x + 10) - 3} = \sqrt{x + 7}$

\vspace{15pt}

Step 2: Find the domain of $\sqrt{x + 7}$: \\
$x + 7 \geq 0$ \\
$x \geq -7$

\vspace{15pt}

The smallest value is $-7$.

\vspace{30pt}

Answer: \colorbox{yellow!100}{-7}
\end{solution}

\vspace{40pt}

% Question 33
% Category: Concept | Difficulty: Medium | Question Type: Multiple Choice
\question
\begin{center}
\begin{tikzpicture}[scale=0.5]
    \draw[->] (-5,0) -- (5,0) node[right] {$x$};
    \draw[->] (0,-1) -- (0,5) node[above] {$y$};
    \draw[ultra thick, blue] (-4,3) -- (4,3);
    \filldraw[blue] (-4,3) circle (2pt);
    \filldraw[blue] (4,3) circle (2pt);
\end{tikzpicture}
\end{center}
What is the range of the function shown in the graph?
\begin{tasks}(4)
    \task $-4 \leq x \leq 4$
    \task $y = 3$
    \task $y = 0$
    \task All real numbers
\end{tasks}

\begin{solution}
\textbf{Conceptual Explanation:} \\
This is a constant function. Regardless of the input $x$ (between -4 and 4), the output $y$ is always exactly 3.

\vspace{30pt}

Answer: \colorbox{yellow!100}{B}
\end{solution}

\vspace{40pt}

% Question 34
% Category: Exam level | Difficulty: Hard | Question Type: Numeric Entry
\question
If the range of $f(x) = x^2 + c$ is $y \geq -5$, what is the value of $c$?

\begin{solution}
\textbf{Calculation and Logic:} \\
The range of a basic $x^2$ function is $y \geq 0$. \\
Adding $c$ shifts the range to $y \geq c$. \\
Since the range is $y \geq -5$, it follows that $c = -5$.

\vspace{30pt}

Answer: \colorbox{yellow!100}{-5}
\end{solution}

\vspace{40pt}

% Question 35
% Category: Practice | Difficulty: Medium | Question Type: Multiple Choice
\question
Which of the following functions has a domain of all real numbers?
\begin{tasks}(2)
    \task $f(x) = \frac{1}{x}$
    \task $f(x) = \sqrt{x}$
    \task $f(x) = 2x + 1$
    \task $f(x) = \frac{1}{x^2}$
\end{tasks}

\begin{solution}
\textbf{Calculation and Logic:} \\
A: Undefined at $x=0$. \\
B: Undefined for $x < 0$. \\
C: Linear functions are defined for all real $x$. \\
D: Undefined at $x=0$.

\vspace{30pt}

Answer: \colorbox{yellow!100}{C}
\end{solution}

\vspace{40pt}

% Question 36
% Category: Exam level | Difficulty: Hard | Question Type: Numeric Entry
\question
$f(x) = (x-2)^2 + 8$ \\
$g(x) = f(x) + 4$ \\
What is the minimum value in the range of function $g$?

\begin{solution}
\textbf{Calculation and Logic:} \\
1. \textbf{Range of f:} The vertex of $f$ is $(2, 8)$. Since it opens up, the range is $y \geq 8$. \\
2. \textbf{Transformation:} $g(x) = f(x) + 4$ shifts the entire graph (and the range) up by 4 units. \\
3. \textbf{New Range:} $y \geq 8 + 4 \Rightarrow y \geq 12$. \\
Minimum value = 12.

\vspace{30pt}

Answer: \colorbox{yellow!100}{12}
\end{solution}

\vspace{40pt}

% Question 37
% Category: Concept | Difficulty: Medium | Question Type: Multiple Choice
\question
\begin{center}
\begin{tikzpicture}[scale=0.6]
    \draw[->] (-1,0) -- (6,0) node[right] {$x$};
    \draw[->] (0,-1) -- (0,5) node[above] {$y$};
    \draw[ultra thick, blue, ->] (1,1) -- (5,4);
    \filldraw[blue] (1,1) circle (2pt);
\end{tikzpicture}
\end{center}
The arrow on the graph indicates the function continues forever in that direction. What is the domain of this function?
\begin{tasks}(2)
    \task $x \geq 1$
    \task $y \geq 1$
    \task $x \geq 0$
    \task All real numbers
\end{tasks}

\begin{solution}
\textbf{Calculation and Logic:} \\
The graph begins horizontally at $x=1$ and the arrow points to the right, covering all values greater than 1.

\vspace{30pt}

Answer: \colorbox{yellow!100}{A}
\end{solution}

\vspace{40pt}

% Question 38
% Category: Exam level | Difficulty: Hard | Question Type: Numeric Entry
\question
If $f(x) = \frac{1}{x - \sqrt{25}}$, for what value of $x$ is the function undefined?

\begin{solution}
\textbf{Calculation and Logic:} \\
$x - \sqrt{25} = 0$ \\
$x - 5 = 0$ \\
$x = 5$

\vspace{30pt}

Answer: \colorbox{yellow!100}{5}
\end{solution}

\vspace{40pt}

% Question 39
% Category: Practice | Difficulty: Hard | Question Type: Multiple Choice
\question
A square has side length $s$. The area $A$ is a function of the side length: $A(s) = s^2$. What is the most appropriate domain for this function in this context?
\begin{tasks}(2)
    \task All real numbers
    \task $s \geq 0$
    \task $s > 0$
    \task $s$ must be an integer
\end{tasks}

\begin{solution}
\textbf{Conceptual Explanation:} \\
In geometry, a side length must be positive to form a shape. While $s=0$ is mathematically possible in the equation, a square with side 0 does not exist. However, the SAT often uses $s > 0$ or $s \geq 0$ for such contexts. Since a square must have "length," $s > 0$ is the best fit.

\vspace{30pt}

Answer: \colorbox{yellow!100}{C}
\end{solution}

\vspace{40pt}

% Question 40
% Category: Exam level | Difficulty: Hard | Question Type: Numeric Entry
\question
$f(x) = \sqrt{x+10}$ \\
If the domain is restricted to $x \leq 26$, what is the maximum value in the range?

\begin{solution}
\textbf{Calculation and Logic:} \\
For an increasing function like $\sqrt{x+c}$, the maximum output occurs at the maximum input. \\
$f(26) = \sqrt{26 + 10} = \sqrt{36} = 6$.

\vspace{30pt}

Answer: \colorbox{yellow!100}{6}
\end{solution}

\vspace{40pt}

% Question 41
% Category: Concept | Difficulty: Medium | Question Type: Multiple Choice
\question
If $f(x) = 5$, what is the range of $f$?
\begin{tasks}(4)
    \task $\{5\}$
    \task All real numbers
    \task $y \geq 5$
    \task $x = 5$
\end{tasks}

\begin{solution}
\textbf{Calculation and Logic:} \\
The only possible output is 5.

\vspace{30pt}

Answer: \colorbox{yellow!100}{A}
\end{solution}

\vspace{40pt}

% Question 42
% Category: Exam level | Difficulty: Hard | Question Type: Numeric Entry
\question
$f(x) = \frac{1}{x^2 + 1}$ \\
What is the largest value in the range of function $f$?

\begin{solution}
\textbf{Conceptual Explanation:} \\
To maximize a fraction with a constant numerator, we must minimize the denominator.

\vspace{30pt}

\textbf{Calculation and Logic:} \\
The smallest value of $x^2 + 1$ is $0 + 1 = 1$ (when $x=0$). \\
The largest value of $f(x)$ is $10 / 1 = 1$. \textbf{(Wait, re-checking numerator...)} \\
$f(0) = 1/1 = 1$. \\
As $x$ gets larger, the denominator gets larger, making the fraction smaller. Thus, 1 is the maximum.

\vspace{30pt}

Answer: \colorbox{yellow!100}{1}
\end{solution}

\vspace{40pt}

% Question 43
% Category: Practice | Difficulty: Medium | Question Type: Multiple Choice
\question
\begin{center}
\begin{tikzpicture}[scale=0.6]
    \draw[->] (-5,0) -- (1,0) node[right] {$x$};
    \draw[->] (0,-1) -- (0,5) node[above] {$y$};
    \draw[ultra thick, blue] (-4,1) -- (-1,4);
\end{tikzpicture}
\end{center}
What is the domain of the segment shown?
\begin{tasks}(2)
    \task $-4 \leq x \leq -1$
    \task $1 \leq y \leq 4$
    \task $-4 \leq y \leq -1$
    \task $x \leq 0$
\end{tasks}

\begin{solution}
\textbf{Calculation and Logic:} \\
Domain is horizontal. The segment spans from $x=-4$ to $x=-1$.

\vspace{30pt}

Answer: \colorbox{yellow!100}{A}
\end{solution}

\vspace{40pt}

% Question 44
% Category: Exam level | Difficulty: Hard | Question Type: Numeric Entry
\question
If $f(x) = \sqrt{kx - 12}$ has a domain of $x \geq 3$, what is the value of $k$?

\begin{solution}
\textbf{Calculation and Logic:} \\
$kx - 12 \geq 0 \Rightarrow kx \geq 12 \Rightarrow x \geq 12/k$. \\
Set $12/k = 3$. \\
$12 = 3k \Rightarrow k = 4$.

\vspace{30pt}

Answer: \colorbox{yellow!100}{4}
\end{solution}

\vspace{40pt}

% Question 45
% Category: Practice | Difficulty: Medium | Question Type: Multiple Choice
\question
Which value of $x$ is in the domain of $f(x) = \sqrt{x-20}$?
\begin{tasks}(4)
    \task 0
    \task 10
    \task 19
    \task 21
\end{tasks}

\begin{solution}
\textbf{Calculation and Logic:} \\
$x - 20 \geq 0 \Rightarrow x \geq 20$. \\
Only 21 is $\geq 20$.

\vspace{30pt}

Answer: \colorbox{yellow!100}{D}
\end{solution}

\vspace{40pt}

% Question 46
% Category: Exam level | Difficulty: Hard | Question Type: Numeric Entry
\question
$f(x) = |x| - 10$ \\
What is the minimum value in the range of $f$?

\begin{solution}
\textbf{Calculation and Logic:} \\
Min of $|x|$ is 0. \\
Min of $0 - 10 = -10$.

\vspace{30pt}

Answer: \colorbox{yellow!100}{-10}
\end{solution}

\vspace{40pt}

% Question 47
% Category: Concept | Difficulty: Medium | Question Type: Multiple Choice
\question
If $f(x) = \sqrt{x}$, how does the domain change if $f(x)$ is transformed to $g(x) = \sqrt{x} + 5$?
\begin{tasks}(1)
    \task The domain increases by 5.
    \task The domain stays the same.
    \task The domain is restricted to $x \geq 5$.
    \task The domain becomes all real numbers.
\end{tasks}

\begin{solution}
\textbf{Conceptual Explanation:} \\
A vertical shift ($+5$ outside the radical) affects the range, not the domain. The input requirement $x \geq 0$ remains the same.

\vspace{30pt}

Answer: \colorbox{yellow!100}{B}
\end{solution}

\vspace{40pt}

% Question 48
% Category: Exam level | Difficulty: Hard | Question Type: Numeric Entry
\question
$f(x) = \frac{1}{x^2 - 10x + 25}$ \\
What value of $x$ is excluded from the domain?

\begin{solution}
\textbf{Calculation and Logic:} \\
$x^2 - 10x + 25 = (x-5)^2 = 0 \Rightarrow x = 5$.

\vspace{30pt}

Answer: \colorbox{yellow!100}{5}
\end{solution}

\vspace{40pt}

% Question 49
% Category: Practice | Difficulty: Medium | Question Type: Multiple Choice
\question
If $f(x) = -x^2$, what is the range?
\begin{tasks}(4)
    \task $y \geq 0$
    \task $y \leq 0$
    \task All real numbers
    \task $x \leq 0$
\end{tasks}

\begin{solution}
\textbf{Calculation and Logic:} \\
$x^2$ is always $\geq 0$. Negating it makes all results $\leq 0$.

\vspace{30pt}

Answer: \colorbox{yellow!100}{B}
\end{solution}

\vspace{40pt}

% Question 50
% Category: Exam level | Difficulty: Hard | Question Type: Numeric Entry
\question
$f(x) = \sqrt{x-5}$ \\
If $f(x) = 10$, what is the value of $x$ in the domain?

\begin{solution}
\textbf{Calculation and Logic:} \\
$10 = \sqrt{x-5} \Rightarrow 100 = x-5 \Rightarrow x = 105$.

\vspace{30pt}

Answer: \colorbox{yellow!100}{105}
\end{solution}

% =============================================================
% Question End
% =============================================================

\end{questions}
\begin{center}
\rule{.5\textwidth}{1pt}
\end{center}
\end{document}
